\documentclass[a4paper,12pt]{article}

\usepackage[a4paper,margin=2.5cm]{geometry}

\input{../preamble}

\newenvironment{referee}[1]
{\noindent {\bf Referee #1}}
{\vspace{1em}}

\newenvironment{response}
{\noindent \color{blue!40!black}{\bf Response}}
{\vspace{1em}}

\bibliographystyle{ieeetr}

\begin{document}

\subsection*{Response to Referee A}
\begin{referee}{A}
	The authors propose the variational method for the thermal states defined by tensor and neural networks.
	It is based on the modified free energy with Renyi entropy, which can be efficiently estimated from Monte Carlo samplings.
	This approach allows us to represent thermal states without an effective imaginary time evolution from high to low temperatures.

	The authors examine the performance of the wider class of a variational ansatz of thermal states compared to the previous study [68] in which MPS representation was checked.
	They found that MPS performs well for the one-dimensional quantum Ising model as the bond dimension increases.
	However, RBM does not perform well.
	For the two-dimensional case, Snakelike SBS ansatz's results improve with increasing the number of strings and the bond dimension, and it is slightly better than MPS, but RBM again does not.
	The computational complexity of the snakelike SBS is scalable.
	Therefore, the snakelike SBS, which includes MPS, may be suitable for the variational ansatz of the thermal states for one and two-dimensional quantum cases.

	Estimating the thermal states of quantum many-body systems is a difficult problem.
	In [68], a new approach using modified free energy has been proposed.
	The authors examined this approach using a wider class of variational ansatz with tensor and neural networks.
	Their results showed that this approach could be more innovative.
	However, it has yet to be.
	In addition, the test examples are limited to the quantum Ising models.
	Therefore, I do not recommend the publication of this manuscript in Physical Review Letters.
\end{referee}

\begin{response}
	We wish to thank the Referee for carefully reading the manuscript.
	We believe that the purpose of this work is to open new avenues for the numerical exploration of the properties of thermal states of many-body systems methods.
	The numerical results are meant to validate and benchmark the method, rather than provide novel physical insights.
	We would also thank them for their detailed feedback and answer to all points below.
\end{response}

\begin{referee}{A}
	In the following, I summarize unclear points from this manuscript:\\
	(1) The concept of MPDO only appears in the caption of Fig.~1.
	It might be better to introduce it in the main text.
	In addition, the following sentence is also written only in the caption: ``All the ansatz used in this work are defined and depicted by their purification.''
\end{referee}

\begin{response}
	[1-2 sentences in main text]
\end{response}

\begin{referee}{A}
	(2) What is the definition of $E$ in Eq.~(3a)?
\end{referee}

\begin{response}
	We thank the Referee for pointing this out.

		[add definition in main text]
\end{response}

\begin{referee}{A}
	(3) Is the definition of $X_i(s)$ correct in the third paragraph of the left column on Page 3?
	In addition, understanding the matrix product state representation through the current description takes much work.
\end{referee}

\begin{response}
	[add 1-2 sentences to clarify]
\end{response}

\begin{referee}{A}
	(4) In Fig.~2(b) and (c)'s caption, $N = 10 \times 10 \rightarrow N = 100$ because the system is one-dimensional.
	What are the MPO results representing the Gibbs ensemble?
\end{referee}

\begin{response}
	We thank the Referee for spotting this inconsistency.
	The corresponding caption has been updated.
\end{response}

\begin{referee}{A}
	(5) I need clarification about the description in Fig.~3's caption for the parameter setting of $h_x$.
\end{referee}

\begin{response}
	The corresponding caption was indeed unclear, it has been updated.
\end{response}

\begin{referee}{A}
	(6) The more additional description is needed for the two-dimensional case.
	Does the $N=16$ mean the square lattice of $4 \times 4$ with an open boundary?
	What type of strings are used?
\end{referee}

\begin{response}
	[Discuss size] \\
	Regarding the string types used in SnakeSBS, we employ strings that snake horizontally and vertically across the lattice, as illustrated in Fig. S4 of the Supplemental Material.
	There are four possible strings in total, denoted as SnakeSBS($n_s=2$) for one horizontal and one vertical string, and SnakeSBS($n_s=4$) for all four strings.
		[add comment in text?]
\end{response}

\subsection*{Response to Referee B}
\begin{referee}{B}
	The article explores an approach to simulate quantum many-body thermal states via a variational procedure minimizing the free energy.
	To make it numerically tractable, it follows the approach of Ref.~[68] by some of the same authors, replacing the von Neuman entropy term in the free energy with the 2-Renyie entropy.
	The article explores a range of variational analyses including matrix
	product states (already explored to some extent in [68]), entangled plaquette states, string-bond states, restricted Boltzmann machine, etc., where the gradient calculations for optimization follow from Monte Carlo sampling.
	The bulk of the presented results focuses on the convergence (or the lack of systematic convergence in some cases) of those ansatzes to the 2-Raneyi ensemble for a minimal example of a paradigmatic Ising model in a transverse field.

	The article is well-written overall and well-motivated.
	The main text is a little dense.
	However, the appendix offers a pedagogical discussion of the technical aspects of the approach, making an entire work sufficiently self-contained.

	While I appreciate the technical aspects of the work, I don't believe it contains material that would warrant publication in the Physical Review Letter.
	The presented results do not look promising enough to make me believe that the approach would have a significant impact of the kind specified in the PRL acceptance criteria.
	However, the article is explorying an interesting approach and I would recommend its publication in a more specialistic journal, such as Physical Review B, after a small extension of the presented data -- see further comments below.
\end{referee}

\begin{response}
	We wish to thank the Referee for carefully reading the manuscript.
	As mentioned above, the purpose of this work is to open new avenues for the numerical exploration of the properties of thermal states of many-body systems methods.
	It has already been acknowledged in several follow-up works (arXiv:2402.07605, arXiv:2404.04005, arXiv:2405.02158) and a recent review paper (arXiv:2402.09402).
	The specific points raised are addressed below.
\end{response}

\begin{referee}{B}
	The Ising model studied by the authors (in a small one-dimensional example where it can be checked against exact diagonalization, as well as in a small 10 by 10 two-dimensional example) is widely believed to be easy to simulate.
	As such, anything but a stellar performance for an Ising model is a prerequisite to consider using a new numerical approach to more challenging problems.
	However, the presented data do not hint at a stellar performance.

	On one hand, there is an issue of uncontrollable approximation replacing von Neuman entropy with the 2-Renyie entropy (independent from the convergence issues of a particular ansatz to the 2-Renyie ensemble).
	The authors make a bold statement that those two ensembles converge in some limits.
	However, there is little evidence quantifying this error -- in a function of inverse temperature beta, and not other quantities that can hide this aspect; as done by the authors in the main text.
	By the way, it is sometimes not clear to me what the authors consider as ``exact'', i.e., when ``exact'' refers to modified distribution, and when to the thermal one.
	This aspect should be made more clear in figure captions.
\end{referee}

\begin{response}
	[explain convergence, expand main text]
\end{response}

\begin{referee}{B}
	Moreover, the data for a two-dimensional example show only high-temperature results, truncating beta at around 0.6 in Fig.~7 -- with a visual deviation of the expectation values from the exact ones around $\beta=0.6$.
	It is hard to judge to what extent this deviation is a result of 2-Renyie approximation or the lack of convergence of variational ansatz.
	The two-dimensional Ising model has a phase transition at around $\beta_c=0.8$ for $h_x = 2.5$ (in the thermodynamic limit).
	As such, the authors omit the most interesting and likely the most challenging feature of the phase diagram of the two-dimensional Ising model -- and it only makes me think why?
	I believe that including the data for larger beta (also for other observables) is necessary for the reader to be able to judge the performance of the approach.
\end{referee}

\begin{response}
	[update figure]
\end{response}

\begin{referee}{B}
	Finally, a small question.
	The authors motivate the necessity of developing the approach by mentioning Hubbard's model.
	How challenging technically would it be to include fermionic statistics in a SnakeMPS ansatz -- that the authors consider as the best performing one -- where various MPS-s would have to match a convention on a global fermionic order in the lattice?
\end{referee}

\begin{response}
	[add comment of fermionic TNs]
\end{response}

\end{document}
